% !PRINTED 282
\weberpage{282}

\begin{center}
\S.\,72.

\gesperrt{Constitution der Gruppe $L_7$ vom Grade 168.}
\end{center}

Unter den hier gefundenen einfachen Gruppen ist die nächste nach der Ikosa\"edergruppe, die hier als $L_5$ wiederkehrt, die Gruppe $L_7$ vom Grade 168, die, wie wir gesehen haben, eine Octa\"edergruppe als Theiler enthält. Da uns diese merkwürdige Gruppe später noch mehrfach begegnen wird, so wollen wir hier noch etwas näher auf ihren Bau eingehen.

Im \S.\,57 ist die Octa\"edergruppe durch drei Elemente $\chi, \omega, \Theta$ in der Form dargestellt:

\weq{1}{\chi^\lambda \omega^u \Theta^v, \quad \lambda = 0, 1, 2; \ \mu = 0, 1; \ \nu = 0, 1, 2, 3.} % [sic] formula prints Latin u, v; the parameter list prints Greek mu, nu

Zwischen diesen Elementen bestehen die Relationen

\weq{2}{\omega\chi = \chi^2\omega, \quad \Theta\omega = \omega\Theta^3, \quad \Theta\chi = \chi^2\omega\Theta^2, \quad \Theta^2\chi = \chi\omega\Theta^3,}

und diese Bedingungen haben wir, wenn noch die Grade 3, 2, 4 der Elemente $\chi, \omega, \Theta$ hinzukommen, als ausreichend nachgewiesen, um das System (1) als Octa\"edergruppe zu charakterisiren.

Aus (2) haben wir im \S.\,57 als Folgerungen die Formeln abgeleitet:

\weq{3}{\begin{gathered}
\omega\chi = \chi^2\omega, \quad \omega\chi^2 = \chi\omega, \quad \Theta^2\chi^2 = \chi^2\omega\Theta,\\
\Theta\chi = \chi^2\omega\Theta^2, \quad \Theta^2\chi = \chi\omega\Theta^3, \quad \Theta^3\chi = \chi^2\Theta,\\
\Theta\omega = \omega\Theta^3, \quad \Theta^2\omega = \omega\Theta^2, \quad \Theta^3\omega = \omega\Theta, \quad \Theta\chi^2 = \chi\Theta^3,
\end{gathered}}

die wir später mehrfach benutzen werden.

Im \S.\,71 haben wir die Octa\"edergruppe auch als Congruenzgruppe nach dem Modul 7 dargestellt und haben in den dortigen Formeln (11)

\weq{4}{\chi = \begin{pmatrix} 0, & 2 \\ 3, & 1 \end{pmatrix}, \ \omega = \begin{pmatrix} 3, & -1 \\ 3, & -3 \end{pmatrix}, \ \Theta = \begin{pmatrix} 2, & 1 \\ 3, & 2 \end{pmatrix}}

gefunden. Für das Folgende ist es aber, im Interesse einer einfacheren Rechnung, zweckmässig, diese Gruppe noch zu trans-
%CONT>>

\bigskip
\hrule
\bigskip

%<<CONT
\gesperrt{Serret} (Cours d'alg\`ebre supérieure) und von C. \gesperrt{Jordan} (Traité des Substitutions). Vergl. \gesperrt{Weber}, Elliptische Functionen etc., \S.\,84 f. Die vollständige Aufstellung aller Theiler rührt von \gesperrt{Gierster} her (Math. Ann., Bd. XVIII). Ausführliche Behandlung der Congruenzgruppen in \glqq \gesperrt{Klein-Fricke}, Vorlesungen über die Theorie der elliptischen Modulfunctionen\grqq{}, Bd. I, Leipzig 1890.
